\documentclass[12pt, a4paper]{article}
\usepackage[utf8]{inputenc}
\usepackage{amsmath, amssymb, amsthm}
\usepackage{geometry}
\usepackage{enumitem}
\usepackage[colorlinks=true, linkcolor=blue!70!black, urlcolor=blue!70!black]{hyperref}
\usepackage{titlesec}
\usepackage{mathrsfs}
\usepackage{xcolor}
\usepackage{mdframed}
\usepackage{tocloft}
\usepackage{fancyhdr}
\usepackage{booktabs}
\usepackage{array}

\geometry{margin=2.5cm, top=3cm, bottom=3cm}

%----------------------------------------------------------
% Page Style
%----------------------------------------------------------
\pagestyle{fancy}
\fancyhf{}
\rhead{\textcolor{gray!70}{\small MA4271 Differential Geometry}}
\lhead{\textcolor{gray!70}{\small Question Bank}}
\cfoot{\thepage}
\renewcommand{\headrulewidth}{0.4pt}

%----------------------------------------------------------
% Colour Palette
%----------------------------------------------------------
\definecolor{sectionblue}{RGB}{30, 80, 160}
\definecolor{hardred}{RGB}{180, 30, 30}
\definecolor{newgreen}{RGB}{0, 110, 50}
\definecolor{newbg}{RGB}{235, 255, 240}
\definecolor{boxbg}{RGB}{240, 245, 255}
\definecolor{rulecol}{RGB}{30, 80, 160}

%----------------------------------------------------------
% Section Formatting
%----------------------------------------------------------
\titleformat{\section}[block]
  {\large\bfseries\color{sectionblue}}
  {\thesection.}{0.5em}{}
  [{\color{rulecol}\titlerule[1.2pt]}]

\titlespacing{\section}{0pt}{18pt}{10pt}

%----------------------------------------------------------
% Question Environment
%----------------------------------------------------------
\newtheoremstyle{qstyle}{4pt}{4pt}{}{}{\bfseries}{.}{.5em}{}
\theoremstyle{qstyle}
\newtheorem{question}{Q}[section]
\renewcommand{\thequestion}{\thesection.\arabic{question}}

% Difficulty badge for hard questions
\newcommand{\hard}{\textcolor{hardred}{\textbf{[$\bigstar$]}}}
% Badge for new questions
\newcommand{\NEW}{\textcolor{newgreen}{\textbf{[NEW]}}}

% Framed box for hard questions
\newenvironment{hardq}[1][]{%
  \begin{mdframed}[backgroundcolor=boxbg, linecolor=hardred, linewidth=1pt,
    innertopmargin=4pt, innerbottommargin=4pt, skipabove=4pt, skipbelow=4pt]
  \begin{question}#1
}{%
  \end{question}
  \end{mdframed}
}

% Framed box for new questions
\newenvironment{newq}[1][]{%
  \begin{mdframed}[backgroundcolor=newbg, linecolor=newgreen, linewidth=1pt,
    innertopmargin=4pt, innerbottommargin=4pt, skipabove=4pt, skipbelow=4pt]
  \begin{question}#1
}{%
  \end{question}
  \end{mdframed}
}

%----------------------------------------------------------
% Math Macros
%----------------------------------------------------------
\newcommand{\R}{\mathbb{R}}
\newcommand{\C}{\mathbb{C}}
\newcommand{\norm}[1]{\left\|#1\right\|}
\newcommand{\abs}[1]{\left|#1\right|}
\newcommand{\ip}[2]{\langle #1,\, #2 \rangle}
\newcommand{\grad}{\nabla}
\DeclareMathOperator{\rank}{rank}
\newcommand{\dd}[2]{\frac{d#1}{d#2}}
\newcommand{\pd}[2]{\frac{\partial #1}{\partial #2}}
\newcommand{\mbf}[1]{\mathbf{#1}}
\newcommand{\kp}{\kappa}
\newcommand{\tp}{\tau}

%----------------------------------------------------------
% TOC Style
%----------------------------------------------------------
\renewcommand{\cftsecfont}{\bfseries\color{sectionblue}}
\renewcommand{\cftsecpagefont}{\bfseries\color{sectionblue}}
\renewcommand{\cftsecdotsep}{\cftdotsep}
\setlength{\cftbeforesecskip}{4pt}

%===========================================================
\begin{document}
%===========================================================

%----------------------------------------------------------
% Title Block
%----------------------------------------------------------
\begin{center}
  {\LARGE \textbf{MA4271 Differential Geometry}}\\[8pt]
  {\large Full Exam-Level Question Bank --- Extended Edition}\\[4pt]
  {\normalsize 40 Questions Per Section $\cdot$ 5 Sections $\cdot$ 200 Questions Total}\\[4pt]
  {\small \textcolor{newgreen}{\textbf{Questions 21--40 in each section are new, drawn directly from lecture notes.}}}\\[6pt]
  \textcolor{gray}{\rule{0.6\linewidth}{0.4pt}}
\end{center}

\vspace{0.5em}

\begin{mdframed}[backgroundcolor=boxbg, linecolor=sectionblue, linewidth=1pt,
  innertopmargin=8pt, innerbottommargin=8pt]
\noindent\textbf{Notes.}\quad All curves are assumed smooth and regular.
Surfaces are in $\R^3$ unless stated otherwise.
Questions marked \hard\ are harder and may require combining ideas from multiple topics.
Questions marked \NEW\ are new additions directly grounded in the NUS MA4271 lecture notes.
Answers should be written with full justification unless stated otherwise.
\end{mdframed}

\vspace{1.5em}

\tableofcontents
\newpage

%===========================================================
\section{Curves and Frenet--Serret Theory}
%===========================================================

\begin{question}
Let $\alpha(t) = (a\cos t,\, a\sin t,\, bt)$ be a circular helix with $a,b>0$.
Reparametrise $\alpha$ by arc length $s$ and compute the Frenet trihedron $\{\mbf{t}, \mbf{n}, \mbf{b}\}$.
Show that $\kp$ and $\tp$ are both constant and express them in terms of $a$ and $b$.
\end{question}

\begin{question}
A curve $\alpha(s)$ parametrised by arc length is a plane curve if and only if $\tau(s) \equiv 0$.
Prove both directions of this equivalence.
\end{question}

\begin{question}
For a non-unit speed curve $\alpha(t)$, show that
\[
  \mbf{t} = \frac{\alpha'}{|\alpha'|}, \qquad
  \mbf{b} = \frac{\alpha' \times \alpha''}{|\alpha' \times \alpha''|}, \qquad
  \kp = \frac{|\alpha' \times \alpha''|}{|\alpha'|^3}.
\]
\end{question}

\begin{question}
Compute the curvature $\kp(t)$ and torsion $\tp(t)$ for the \textbf{Catenary Helix}
$\alpha(t) = (\cosh t,\, \sinh t,\, t)$.
\end{question}

\begin{question}
For $\alpha(t) = (t, t^2, t^3)$, find the Frenet trihedron $\{\mbf{t}, \mbf{n}, \mbf{b}\}$ at $t=1$.
Compute $\kp(1)$ and $\tp(1)$.
\end{question}

\begin{hardq}
  \textbf{Lancret's Theorem.} \hard\;
  Prove that a curve is a \emph{generalized helix} (makes a constant angle with a fixed direction) if and only if $\kp/\tp$ is constant along the curve.
\end{hardq}

\begin{question}
Let $\alpha(s)$ be unit-speed. Derive the Taylor expansion
\[
  |\alpha(s)-\alpha(0)|^2 = s^2 - \tfrac{\kp(0)^2}{12}s^4 + O(s^5)
\]
and hence obtain the leading-order formula $|\alpha(s)-\alpha(0)| = s - \frac{\kp(0)^2}{24}s^3 + O(s^4)$.
\end{question}

\begin{question}
Let $\alpha(s)$ be a unit-speed curve lying on a sphere of radius $R$.
Deduce the constraint
\[
  R^2 = \frac{1}{\kp^2} + \left(\frac{1}{\tp}\cdot\frac{d}{ds}\frac{1}{\kp}\right)^2.
\]
\end{question}

\begin{question}
Find the curvature and torsion of
$\alpha(t) = \bigl(t - \sin t,\; 1 - \cos t,\; 4\sin(t/2)\bigr)$.
\end{question}

\begin{hardq}
  \textbf{Darboux Vector.} \hard\;
  Show that the Frenet equations $\mbf{t}' = \kp\mbf{n}$, $\mbf{n}' = -\kp\mbf{t}+\tp\mbf{b}$, $\mbf{b}'=-\tp\mbf{n}$
  can be written uniformly as $\mbf{v}' = \mbf{d}\times\mbf{v}$ for $\mbf{v}\in\{\mbf{t},\mbf{n},\mbf{b}\}$,
  where $\mbf{d}=\tp\mbf{t}+\kp\mbf{b}$ is the \emph{Darboux vector}.
\end{hardq}

\begin{question}
\textbf{Evolute and Involute.}\;
The \emph{evolute} of a unit-speed plane curve $\alpha(s)$ is $\beta(s) = \alpha(s) + \frac{1}{\kp(s)}\mbf{n}(s)$.
Show that $\beta'(s)$ is parallel to $\mbf{b}(s)$ and interpret $\beta$ geometrically.
\end{question}

\begin{question}
\textbf{Fenchel's Theorem.}\;
State Fenchel's Theorem and apply it to deduce that a closed curve with $\kp < 2\pi/L$ everywhere cannot be convex.
\end{question}

\begin{question}
\textbf{Bertrand Curves.}\;
Show that $\alpha$ and $\beta$ are Bertrand mates if and only if there exist constants $a, b$ such that $a\kp + b\tp = 1$.
\end{question}

\begin{question}
For $\alpha(t) = (e^t \cos t,\; e^t \sin t,\; e^t)$, compute the arc length from $t=0$ to $t=T$, find $\mbf{t}(t)$, and show $\kp(t)$ is constant.
\end{question}

\begin{question}
\textbf{Four-Vertex Theorem.}\;
State the theorem and verify it for an ellipse $\alpha(t) = (a\cos t, b\sin t)$ with $a\neq b$ by computing where $\kp'(t)=0$.
\end{question}

\begin{question}
Let $\alpha(s)$ be unit-speed. Prove the Frenet--Serret equations by differentiating $\mbf{t}\cdot\mbf{t}=1$, $\mbf{t}\cdot\mbf{n}=0$, etc., from first principles.
\end{question}

\begin{hardq}
  \textbf{Isoperimetric Inequality.} \hard\;
  State $L^2 \geq 4\pi A$. Sketch a Fourier series proof and characterise equality.
\end{hardq}

\begin{question}
\textbf{Signed Curvature.}\;
For $\alpha(t)=(x(t),y(t))$, define $\kp_s = (x'y''-y'x'')/|\alpha'|^3$. Show it changes sign under orientation reversal and $|\kp_s|=\kp$.
\end{question}

\begin{question}
\textbf{Umlaufsatz.}\;
State Hopf's theorem and verify total signed curvature equals $\pm 2\pi$ for (a) a circle of radius $r$, (b) an ellipse.
\end{question}

\begin{hardq}
  \textbf{Fundamental Theorem of Curve Theory.} \hard\;
  State the theorem and outline the proof via the Frenet ODE system.
\end{hardq}

%---------- NEW QUESTIONS 21-40 ----------
\begin{newq}
  \NEW\; \textbf{Smooth vs.\ Singular.}\;
  The lecture notes define a parametrised curve $\alpha(t)$ to be \emph{non-singular} if $\alpha'(t)\neq\mathbf{0}$ for all $t$.
  \begin{enumerate}[label=(\alph*)]
    \item Let $\alpha(t)=(t^2, t^3)$. Show $\alpha$ is smooth but singular at $t=0$.
    \item Explain geometrically why no velocity vector can reasonably be assigned at the origin.
    \item Identify the value of $|\alpha'(t)|$ near $t=0$ and explain why this prevents arc-length reparametrisation.
  \end{enumerate}
\end{newq}

\begin{newq}
  \NEW\; \textbf{Arc-Length Reparametrisation Step by Step.}\;
  Let $\gamma(t) = (r\cos t, r\sin t, 0)$ for $t\in(0,\pi)$.
  \begin{enumerate}[label=(\alph*)]
    \item Compute the arc-length function $s(t) = \int_0^t |\gamma'(u)|\,du$.
    \item Find the inverse $t = t(s)$ explicitly.
    \item Write $\alpha(s) = \gamma(t(s))$ parametrised by arc length and verify $|\alpha'(s)|=1$.
  \end{enumerate}
  Then explain why the Inverse Function Theorem guarantees the existence and smoothness of $t(s)$ for any non-singular smooth curve.
\end{newq}

\begin{newq}
  \NEW\; \textbf{Unit Vectors and Perpendicularity.}\;
  The lecture notes prove that if $\mathbf{v}(t)$ is a unit vector for all $t$, then $\mathbf{v}'(t)\perp\mathbf{v}(t)$.
  \begin{enumerate}[label=(\alph*)]
    \item Supply the proof by differentiating $\mathbf{v}(t)\cdot\mathbf{v}(t)=1$.
    \item Apply this lemma to deduce that $\alpha''(s)\perp\alpha'(s)$ for any arc-length parametrised curve.
    \item Use it again to show $\mathbf{b}'(s)\perp\mathbf{b}(s)$.
  \end{enumerate}
\end{newq}

\begin{newq}
  \NEW\; \textbf{Plane Curves and the Binormal.}\;
  The lecture notes prove: if $\mathbf{b}(s)=\mathbf{c}$ is constant, then $\alpha(s)$ lies on a plane perpendicular to $\mathbf{c}$.
  \begin{enumerate}[label=(\alph*)]
    \item Reproduce the proof: let $z(s) = \alpha(s)\cdot\mathbf{c}$, differentiate, and conclude $z(s)$ is constant.
    \item Verify that the constant equals $\alpha(0)\cdot\mathbf{c}$. Interpret geometrically.
    \item Conversely, if $\alpha(s)$ lies on a plane with unit normal $\mathbf{c}$, prove $\mathbf{b}(s)=\pm\mathbf{c}$ and hence $\tp\equiv 0$.
  \end{enumerate}
\end{newq}

\begin{newq}
  \NEW\; \textbf{Deriving $\mathbf{n}'(s)$.}\;
  The lecture notes derive $\mathbf{n}'(s) = -\kp(s)\mathbf{t}(s) - \tp(s)\mathbf{b}(s)$ by differentiating $\mathbf{n}=\mathbf{b}\wedge\mathbf{t}$.
  \begin{enumerate}[label=(\alph*)]
    \item Carry out the derivation step by step, applying the product rule for the cross product.
    \item Write the Frenet--Serret system as the matrix equation
    $\begin{pmatrix}\mathbf{t}'\\\mathbf{n}'\\\mathbf{b}'\end{pmatrix}
    = \begin{pmatrix}0&\kp&0\\-\kp&0&\tp\\0&-\tp&0\end{pmatrix}
    \begin{pmatrix}\mathbf{t}\\\mathbf{n}\\\mathbf{b}\end{pmatrix}$
    and observe that the matrix is skew-symmetric.
    \item Explain why skew-symmetry follows from $\{\mathbf{t},\mathbf{n},\mathbf{b}\}$ being orthonormal.
  \end{enumerate}
\end{newq}

\begin{newq}
  \NEW\; \textbf{The Worked Example from Section~1.4.}\;
  Consider the arc-length parametrised curve
  \[
    \alpha(s) = \Bigl(\sqrt{2}\arctan\tfrac{s}{\sqrt{2}},\;\log(s^2+2),\;s - \sqrt{2}\arctan\tfrac{s}{\sqrt{2}}\Bigr).
  \]
  \begin{enumerate}[label=(\alph*)]
    \item Verify that $|\alpha'(s)|=1$.
    \item Compute $\kp(s)$ and $\tp(s)$ as in the lecture notes.
    \item Independently verify the formula $\mathbf{n}'(s)=-\kp(s)\mathbf{t}(s)-\tp(s)\mathbf{b}(s)$ for this example by direct differentiation.
  \end{enumerate}
\end{newq}

\begin{newq}
  \NEW\; \textbf{Frenet Frame of the Standard Example in Section~1.5.}\;
  The lecture notes compute the Frenet frame for $\alpha(t)=(3t-t^3, 3t^2, 3t+t^3)$.
  \begin{enumerate}[label=(\alph*)]
    \item Using the non-unit-speed formulas from Section~1.5, reproduce the computation of $\kp(t)$ and $\tp(t)$.
    \item Verify that $\kp = -\tp = \dfrac{1}{3(1+t^2)^2}$.
    \item Compute the Frenet trihedron $\{\mathbf{t},\mathbf{n},\mathbf{b}\}$ at $t=0$ and $t=1$.
  \end{enumerate}
\end{newq}

\begin{newq}
  \NEW\; \textbf{Rigid Motions and the Frenet Frame.}\;
  A rigid motion is $R(\mathbf{x})=A\mathbf{x}+\mathbf{v}$ where $A\in O(3)$ with $\det A=1$ and $\mathbf{v}\in\R^3$.
  \begin{enumerate}[label=(\alph*)]
    \item Show $R$ is an isometry: $|R(\mathbf{x})-R(\mathbf{y})|=|\mathbf{x}-\mathbf{y}|$ for all $\mathbf{x},\mathbf{y}$.
    \item Let $\beta(s)=R(\alpha(s))$. Show $\beta'(s)=A\alpha'(s)$ and hence the unit tangent of $\beta$ equals $A\mathbf{t}_\alpha$.
    \item Conclude that $\kp_\beta(s)=\kp_\alpha(s)$ and $\tp_\beta(s)=\tp_\alpha(s)$.
  \end{enumerate}
\end{newq}

\begin{newq}
  \NEW\; \textbf{Uniqueness in the Fundamental Theorem.}\;
  Suppose $\alpha(s)$ and $\beta(s)$ satisfy the same curvature $\kp(s)$ and torsion $\tp(s)$, and at $s=0$ their Frenet trihedra are aligned: $\mathbf{t}_1(0)=\mathbf{t}_2(0)$, $\mathbf{n}_1(0)=\mathbf{n}_2(0)$, $\mathbf{b}_1(0)=\mathbf{b}_2(0)$.
  \begin{enumerate}[label=(\alph*)]
    \item Define $f(s) = |\mathbf{t}_1(s)-\mathbf{t}_2(s)|^2+|\mathbf{n}_1(s)-\mathbf{n}_2(s)|^2+|\mathbf{b}_1(s)-\mathbf{b}_2(s)|^2$.
    \item Show $f'(s)=0$ using the Frenet equations and conclude $f\equiv 0$.
    \item Integrate $\mathbf{t}_1(s)=\mathbf{t}_2(s)$ to deduce $\alpha(s)=\beta(s)$.
  \end{enumerate}
\end{newq}

\begin{newq}
  \NEW\; \textbf{Non-Singular Condition and Arc Length.}\;
  The lecture notes use the Inverse Function Theorem to reparametrise by arc length.
  \begin{enumerate}[label=(\alph*)]
    \item State the one-variable Inverse Function Theorem as stated in Section~1.1 of the notes.
    \item Let $\alpha(t)$ be non-singular. Show $s(t)=\int_0^t|\alpha'(u)|\,du$ satisfies $s'(t)>0$.
    \item Explain why $s'(t)>0$ implies $t(s)$ exists and is smooth, even when the integral $s(t)$ cannot be evaluated in closed form.
  \end{enumerate}
\end{newq}

\begin{newq}
  \NEW\; \textbf{Translating Curves.}\;
  The lecture notes observe: if $\beta(s)=\alpha(s)+\mathbf{c}$ for a constant vector $\mathbf{c}$, then $\alpha$ and $\beta$ have the same Frenet data.
  \begin{enumerate}[label=(\alph*)]
    \item Prove this from scratch: $\beta'=\alpha'$, so $\mathbf{t}_\beta=\mathbf{t}_\alpha$, and so on.
    \item Explain why translation is a special case of a rigid motion.
    \item Give an example of two curves that share the same $\kp$ and $\tp$ but differ by a rotation (not just a translation).
  \end{enumerate}
\end{newq}

\begin{newq}
  \NEW\; \textbf{Curvature via the Osculating Circle.}\;
  The curvature $\kp(s_0)=1/R$ where $R$ is the radius of the osculating (best-approximating) circle at $\alpha(s_0)$.
  \begin{enumerate}[label=(\alph*)]
    \item For a circle $\alpha(s) = (r\cos(s/r), r\sin(s/r))$, verify $\kp(s)=1/r$ from the definition $\kp=|\alpha''|$.
    \item For a straight line $\alpha(s)=\mathbf{a}+s\mathbf{t}_0$, verify $\kp=0$ and explain why the radius of curvature is infinite.
    \item Using the centripetal acceleration analogy from Section~1.3 of the notes, explain intuitively why higher curvature corresponds to a tighter bend.
  \end{enumerate}
\end{newq}

\begin{newq}
  \NEW\; \textbf{Orientation and the Cross Product.}\;
  Section~1.2 of the notes discusses orientation. Let $\{\mathbf{e}_1,\mathbf{e}_2,\mathbf{e}_3\}$ be a right-handed basis.
  \begin{enumerate}[label=(\alph*)]
    \item Prove the product rule for cross products: $(\mathbf{v}(t)\wedge\mathbf{w}(t))'=\mathbf{v}'(t)\wedge\mathbf{w}(t)+\mathbf{v}(t)\wedge\mathbf{w}'(t)$.
    \item Use this to verify $\mathbf{b}'(s) = ({\mathbf{t}(s)\wedge\mathbf{n}(s)})'= \tp(s)\mathbf{n}(s)$ without assuming the result.
    \item What happens to the orientation of the Frenet frame under a reflection? Does $\kp$ or $\tp$ change sign?
  \end{enumerate}
\end{newq}

\begin{newq}
  \NEW\; \textbf{Speed, Velocity, Acceleration.}\;
  The lecture notes carefully distinguish velocity $\alpha'(t)$, speed $|\alpha'(t)|$, and acceleration $\alpha''(t)$.
  \begin{enumerate}[label=(\alph*)]
    \item For a non-arc-length curve $\alpha(t)$, write $\alpha''(t)$ in terms of $\mathbf{t}$, $\mathbf{n}$, $s'(t)=|\alpha'(t)|$, and $\kp$.
    Specifically show $\alpha''(t) = \dfrac{d|\alpha'|}{dt}\mathbf{t} + \kp|\alpha'|^2\mathbf{n}$.
    \item Identify the tangential component of acceleration (change in speed) and the normal component (centripetal term).
    \item Explain the \emph{Warning} in Section~1.3: why does using acceleration to gauge curvature fail for non-arc-length curves?
  \end{enumerate}
\end{newq}

\begin{newq}
  \NEW\; \textbf{Frenet Frame of the $\alpha(t)=(t,t^2,t^3)$ Curve Revisited.}\;
  Recompute the Frenet frame for $\alpha(t)=(t,t^2,t^3)$ using the non-unit-speed formulas of Section~1.5.
  \begin{enumerate}[label=(\alph*)]
    \item Compute $\alpha'$, $\alpha''$, $\alpha'''$.
    \item Use $\mathbf{b}=\dfrac{\alpha'\times\alpha''}{|\alpha'\times\alpha''|}$ to find $\mathbf{b}(t)$ for general $t$.
    \item Verify: as $t\to 0$, the torsion $\tp(t)\to 2/3$ and the curvature $\kp(0)=2$.
  \end{enumerate}
\end{newq}

\begin{newq}
  \NEW\; \textbf{Elliptical Orbit.}\;
  A particle in an elliptical orbit (in the $xy$-plane) follows $\alpha(t)=(a\cos t, b\sin t, 0)$ with $a>b>0$.
  \begin{enumerate}[label=(\alph*)]
    \item Compute $\kp(t)$ using the formula for non-unit-speed curves.
    \item Show $\kp$ is maximum at $t=0,\pi$ (ends of the major axis) and minimum at $t=\pi/2, 3\pi/2$.
    \item Using the Frenet frame, express the curvature at $t=0$ and $t=\pi/2$ in terms of $a$ and $b$.
  \end{enumerate}
\end{newq}

\begin{newq}
  \NEW\; \textbf{Helical Curve Check.}\;
  A curve $\alpha(s)$ is called a \emph{circular helix} if $\kp=\kp_0>0$ and $\tp=\tp_0$ are both constant.
  \begin{enumerate}[label=(\alph*)]
    \item Starting from the Frenet ODEs with constant $\kp_0, \tp_0$, reduce the system to a second-order ODE for $\mathbf{t}(s)$.
    \item Solve the ODE to show $\mathbf{t}(s)$ traces a circle on the unit sphere.
    \item Integrate $\mathbf{t}(s)$ to recover $\alpha(s)$ and verify it is a helix with the given $a,b$ (in terms of $\kp_0, \tp_0$).
  \end{enumerate}
\end{newq}

\begin{newq}
  \NEW\; \textbf{Product Rule Application.}\;
  Suppose $\alpha(s)$ is unit-speed and define $f(s) = \alpha(s)\cdot\mathbf{c}$ for a fixed unit vector $\mathbf{c}$.
  \begin{enumerate}[label=(\alph*)]
    \item Compute $f'(s)$ and $f''(s)$.
    \item Show that if $\mathbf{c}$ is chosen to be the binormal $\mathbf{b}(0)$ at $s=0$, then $f'(0)=0$ and $f''(0)=-\kp(0)\langle\mathbf{n}(0),\mathbf{b}(0)\rangle=0$.
    \item Using this, explain why the curve remains close to its osculating plane for small $s$: $\alpha(s)\cdot\mathbf{b}(0) = O(s^3)$.
  \end{enumerate}
\end{newq}

\begin{newq}
  \NEW\; \textbf{Torsion as a Measure of Non-Planarity.}\;
  \begin{enumerate}[label=(\alph*)]
    \item Suppose $\tp(s_0)>0$. Give an intuitive and a rigorous explanation of what this says about how the curve is twisting out of its osculating plane near $s_0$.
    \item If $\kp>0$ and $\tp\equiv 0$, prove the curve lies in a fixed plane by showing $\mathbf{b}$ is constant.
    \item If $\kp\equiv 0$, show $\alpha(s)$ is a straight line regardless of $\tp$.
  \end{enumerate}
\end{newq}

\begin{newq}
  \NEW\; \textbf{Creative: The Train on the Track.}\;
  A passenger on a train experiences forces that tell them how the track curves.
  The train moves along $\alpha(t)$ with non-constant speed.
  \begin{enumerate}[label=(\alph*)]
    \item Decompose $\alpha''(t)$ into components along $\mathbf{t}$ and $\mathbf{n}$ and identify which component the passenger feels as ``sideways force'' (centripetal).
    \item Suppose the train moves at constant speed $v_0$. Show the sideways force is $m\kp v_0^2$, proportional to curvature.
    \item The train enters a spiral easement where $\kp(t)=ct$ (linearly increasing). How does the sideways force change over time? Why are spiral easements used in real railway design?
  \end{enumerate}
\end{newq}

\begin{newq}
  \NEW\; \textbf{Creative: Knottedness and Total Torsion.}\;
  For a closed curve $\alpha: [0,L]\to\R^3$ with $\alpha(0)=\alpha(L)$:
  \begin{enumerate}[label=(\alph*)]
    \item The total torsion is $\int_0^L \tp(s)\,ds$. Show this equals the total rotation of the binormal vector around the curve.
    \item For the standard unknot (a circle), compute the total torsion.
    \item For the $(2,1)$-torus knot $\alpha(t)=((R+r\cos 2t)\cos t, (R+r\cos 2t)\sin t, r\sin 2t)$, describe qualitatively why the total torsion is nonzero. (No explicit computation required; argue geometrically.)
  \end{enumerate}
\end{newq}

%===========================================================
\section{Regular Surfaces and Parametrisations}
%===========================================================

\begin{question}
State the \textbf{Pre-image Theorem}.
Show that the ellipsoid $\frac{x^2}{a^2} + \frac{y^2}{b^2} + \frac{z^2}{c^2} = 1$ is a regular surface.
\end{question}

\begin{question}
Is the cone $z^2 = x^2 + y^2$ a regular surface?
Examine $\grad f$ at $(0,0,0)$ and explain the geometric failure.
\end{question}

\begin{question}
Parametrise a surface of revolution. State necessary and sufficient conditions on $f$ and $g$ for regularity.
\end{question}

\begin{question}
Show that stereographic projection of $S^2\setminus\{N\}$ onto $\R^2$ is a diffeomorphism.
\end{question}

\begin{question}
Let $S = \{f(x,y,z)=c\}$ with $\grad f\neq 0$ on $S$. Show $\grad f(p)$ is normal to $T_pS$.
\end{question}

\begin{question}
Compute $E, F, G$ for the catenoid $\mathbf{x}(u,v) = (c\cosh v\cos u,\; c\cosh v\sin u,\; cv)$.
\end{question}

\begin{question}
Compute the surface area of the torus $\mathbf{x}(u,v) = \bigl((R+r\cos u)\cos v,\; (R+r\cos u)\sin v,\; r\sin u\bigr)$.
\end{question}

\begin{question}
Show that the graph $\{(x,y,f(x,y))\}$ of any smooth $f:\R^2\to\R$ is always a regular surface.
\end{question}

\begin{question}
Show $\mathbf{x}(u,v)=(u+v,\;u-v,\;uv)$ defines a regular surface and compute $E,F,G$.
\end{question}

\begin{hardq}
  \textbf{Orthogonal Parametrisations.} \hard\;
  Explain why locally one can find an orthogonal parametrisation with $F=0$.
\end{hardq}

\begin{question}
\textbf{Transition Maps.}\;
Show the transition map $\mathbf{y}^{-1}\circ\mathbf{x}$ between overlapping parametrisations is a diffeomorphism.
\end{question}

\begin{question}
\textbf{Tangent Space as a Subspace.}\;
Define $T_pS$ intrinsically and show it is a 2-dimensional vector subspace of $\R^3$.
\end{question}

\begin{question}
\textbf{Implicit Function Theorem.}\;
Show that $f(x,y,z)=0$ with $\partial f/\partial z\neq 0$ can locally be written as a graph $z=g(x,y)$.
\end{question}

\begin{question}
Compute $E,F,G$ for the helicoid $\mathbf{x}(u,v)=(v\cos u,\; v\sin u,\; cu)$ and find $\mathbf{N}$.
\end{question}

\begin{question}
Show that the area formula $\iint|\mathbf{x}_u\times\mathbf{x}_v|\,du\,dv$ is independent of parametrisation.
\end{question}

\begin{question}
Parametrise the M\"obius strip and show it is not orientable.
\end{question}

\begin{question}
Define $d\varphi_p: T_pS_1\to T_{\varphi(p)}S_2$ and show it is a well-defined linear map.
\end{question}

\begin{hardq}
  \textbf{Compact Surfaces.} \hard\;
  Show every compact regular surface in $\R^3$ has at least one point of positive Gaussian curvature.
\end{hardq}

\begin{question}
\textbf{Conformal Maps.}\;
Show $\varphi$ is conformal iff $E_1=\lambda^2 E_2$, $F_1=\lambda^2 F_2$, $G_1=\lambda^2 G_2$.
\end{question}

\begin{question}
\textbf{Tubular Neighbourhood.}\;
For compact $S$, show that for small $\varepsilon>0$ every point of $T_\varepsilon$ has a unique closest point on $S$.
\end{question}

%---------- NEW QUESTIONS 21-40 ----------

\begin{newq}
  \NEW\; \textbf{The Three Conditions for a Regular Surface.}\;
  The lecture notes define a regular surface using three conditions on a coordinate function $\mathbf{x}: U\to S$.
  \begin{enumerate}[label=(\alph*)]
    \item State all three conditions precisely.
    \item For the unit sphere, exhibit the coordinate function covering the northern hemisphere and verify each condition explicitly, including the RREF computation of $d\mathbf{x}_q$.
    \item Explain why the six coordinate functions are needed to cover the full sphere.
  \end{enumerate}
\end{newq}

\begin{newq}
  \NEW\; \textbf{The Union of Two Coordinate Planes.}\;
  Let $S=\{(x,y,z)\in\R^3 : yz=0\}$ (the union of the $xy$- and $xz$-planes).
  \begin{enumerate}[label=(\alph*)]
    \item Show that every point other than $(x,0,0)$ has a coordinate function by checking the three conditions.
    \item Explain geometrically why no coordinate function can cover a neighbourhood of $(1,0,0)$ in $S$.
    \item Relate this to the fact that $S$ is not a topological manifold at such points.
  \end{enumerate}
\end{newq}

\begin{newq}
  \NEW\; \textbf{Surface of Revolution: Full Verification.}\;
  The lecture notes treat the surface of revolution $\mathbf{x}(u,v)=(f(v)\cos u, f(v)\sin u, v)$ for $f(v)>0$.
  \begin{enumerate}[label=(\alph*)]
    \item Verify the three conditions for $\mathbf{x}: (0,2\pi)\times\R \to S\cap V$ where $V=\R^3\setminus\{(x,0,z):x\geq 0\}$.
    \item Why is a second coordinate function needed? Provide it explicitly.
    \item Show the transition map between the two coordinate functions is smooth.
  \end{enumerate}
\end{newq}

\begin{newq}
  \NEW\; \textbf{Regularity via the Inverse Function Theorem.}\;
  Section~2.4 uses the Inverse Function Theorem to show level sets $f^{-1}(a)$ are regular surfaces when $\grad f\neq\mathbf{0}$.
  \begin{enumerate}[label=(\alph*)]
    \item State the multivariable Inverse Function Theorem as in Section~2.3.
    \item Apply it to show that the hyperboloid $\{z^2=x^2+y^2+1\}$ is a regular surface.
    \item Compare the efficiency of this approach with directly exhibiting coordinate functions.
  \end{enumerate}
\end{newq}

\begin{newq}
  \NEW\; \textbf{Tangent Vectors from Curves.}\;
  Section~2.8 defines $T_pS$ as the set of velocity vectors of smooth curves through $p$.
  \begin{enumerate}[label=(\alph*)]
    \item Let $\mathbf{x}(u,v)=(u,v,\sqrt{1-u^2-v^2})$ (northern hemisphere) and $p=\mathbf{x}(\tfrac{1}{2},\tfrac{1}{3})$.
    Compute $\mathbf{x}_u(p)$ and $\mathbf{x}_v(p)$.
    \item Find a curve $\alpha(t)$ on the sphere such that $\alpha(0)=p$ and $\alpha'(0) = 2\mathbf{x}_u(p)+3\mathbf{x}_v(p)$, following the recipe in the notes.
    \item Explain why $T_pS$ is the column space of $d\mathbf{x}_p$ and is independent of the choice of coordinate function.
  \end{enumerate}
\end{newq}

\begin{newq}
  \NEW\; \textbf{Smooth Functions on Surfaces.}\;
  Let $f(x,y,z)=xz^2$ on the unit sphere.
  \begin{enumerate}[label=(\alph*)]
    \item Using $\mathbf{x}(u,v)=(u,v,\sqrt{1-u^2-v^2})$ for the northern hemisphere, compute $(f\circ\mathbf{x})(u,v)$.
    \item Compute $\dfrac{\partial f}{\partial u}$ and $\dfrac{\partial^2 f}{\partial u^2}$ at the north pole $(u,v)=(0,0)$.
    \item Prove that $f$ is smooth on the sphere by verifying it on each of the six hemisphere patches.
  \end{enumerate}
\end{newq}

\begin{newq}
  \NEW\; \textbf{Independence of Smoothness from Charts.}\;
  Let $S$ be a regular surface, $f:S\to\R$ smooth at $p$ using chart $\mathbf{x}$, and let $\mathbf{y}$ be another chart also covering $p$.
  \begin{enumerate}[label=(\alph*)]
    \item State the transition map $h=\mathbf{x}^{-1}\circ\mathbf{y}$ and explain why it is smooth (quoting Section~2.6).
    \item Show that $(f\circ\mathbf{y}) = (f\circ\mathbf{x})\circ h$, hence $(f\circ\mathbf{y})$ is smooth.
    \item Conclude that the definition of smoothness of $f:S\to\R$ is independent of the coordinate chart.
  \end{enumerate}
\end{newq}

\begin{newq}
  \NEW\; \textbf{Smooth Maps Between Surfaces.}\;
  Let $\varphi: S_1\to S_2$ be the antipodal map $\varphi(x,y,z)=(-x,-y,-z)$ on the unit sphere.
  \begin{enumerate}[label=(\alph*)]
    \item Verify that $\varphi$ is smooth on the northern hemisphere by computing $(\mathbf{x}_2^{-1}\circ\varphi\circ\mathbf{x}_1)(u,v)=(-u,-v)$.
    \item Explain how the short proof from Section~2.10 (using the restriction of a smooth map on $\R^3$) works.
    \item Compute $d\varphi_p: T_pS\to T_{\varphi(p)}S$ explicitly for $p=(\tfrac{1}{2},\tfrac{1}{2},\tfrac{1}{\sqrt{2}})$.
  \end{enumerate}
\end{newq}

\begin{newq}
  \NEW\; \textbf{The Differential of a Smooth Map.}\;
  Section~2.11 defines $d\varphi_p(\mathbf{v}_1) = \beta'(0)$ where $\beta=\varphi\circ\alpha$ and $\alpha'(0)=\mathbf{v}_1$.
  \begin{enumerate}[label=(\alph*)]
    \item State the key lemma expressing $d\varphi$ as $d\mathbf{x}_2\cdot d\Phi\cdot(a,b)^T$ where $\Phi=\mathbf{x}_2^{-1}\circ\varphi\circ\mathbf{x}_1$.
    \item Verify that this is a linear map (i.e.\ the matrix $d\mathbf{x}_2\cdot d\Phi$ acts linearly on $(a,b)^T$).
    \item Apply the lemma to the $90^\circ$ rotation example from the lecture: compute $d\varphi_p$ for the specific point $p=(\frac{3}{13},\frac{4}{13},\frac{12}{13})$.
  \end{enumerate}
\end{newq}

\begin{newq}
  \NEW\; \textbf{First Fundamental Form for the Tube.}\;
  The cylinder (tube) is $\mathbf{x}(u,v)=(\cos u, \sin u, v)$.
  \begin{enumerate}[label=(\alph*)]
    \item Compute $\mathbf{x}_u$ and $\mathbf{x}_v$; hence find $E, F, G$.
    \item Express the arc length of a curve $\alpha(t)=\mathbf{x}(u(t),v(t))$ on the tube in terms of $u'$ and $v'$.
    \item Show that the first fundamental form for the tube is identical to that of a flat plane in Cartesian coordinates.  What does this tell us about the tube and the plane?
  \end{enumerate}
\end{newq}

\begin{newq}
  \NEW\; \textbf{Area Formula via $\sqrt{EG-F^2}$.}\;
  \begin{enumerate}[label=(\alph*)]
    \item Prove the identity $\|\mathbf{x}_u\wedge\mathbf{x}_v\|^2 = EG-F^2$ from the lecture notes.
    \item Use it to re-derive the area of the torus $\mathbf{x}(u,v)=((a+r\cos u)\cos v,(a+r\cos u)\sin v,r\sin u)$, obtaining $A=4\pi^2 ar$.
    \item Explain why $\sqrt{EG-F^2}\,du\,dv$ is the correct area element and relate it to the cross product of the tangent vectors.
  \end{enumerate}
\end{newq}

\begin{newq}
  \NEW\; \textbf{Change of Parameters as a Diffeomorphism.}\;
  Section~2.6 proves that the transition map $h=\mathbf{x}_1^{-1}\circ\mathbf{x}_2: V'\to U'$ is a diffeomorphism.
  \begin{enumerate}[label=(\alph*)]
    \item Show $h$ is a bijection (as in the blog proof).
    \item Using Proposition~11 from Section~2.6, prove $h$ is smooth.
    \item Give an explicit example with the sphere: compute $h$ between the charts $\mathbf{x}_1(u,v)=(u,v,\sqrt{1-u^2-v^2})$ and $\mathbf{x}_2(\xi,\eta)=(\sqrt{1-\xi^2-\eta^2},\xi,\eta)$ on their overlap.
  \end{enumerate}
\end{newq}

\begin{newq}
  \NEW\; \textbf{Angle Between Tangent Vectors.}\;
  Let $\theta$ be the angle between $\mathbf{x}_u$ and $\mathbf{x}_v$.
  \begin{enumerate}[label=(\alph*)]
    \item Show $\cos\theta = F/\!\sqrt{EG}$.
    \item For an orthogonal parametrisation ($F=0$), show $\mathbf{x}_u\perp\mathbf{x}_v$.
    \item Let $\mathbf{v}_1=a\mathbf{x}_u+b\mathbf{x}_v$ and $\mathbf{v}_2=c\mathbf{x}_u+d\mathbf{x}_v$. Derive a formula for the angle between $\mathbf{v}_1$ and $\mathbf{v}_2$ in terms of $E,F,G,a,b,c,d$.
  \end{enumerate}
\end{newq}

\begin{newq}
  \NEW\; \textbf{Regular vs.\ Non-Regular: The Cone.}\;
  \begin{enumerate}[label=(\alph*)]
    \item Show that $f(x,y,z)=x^2+y^2-z^2$ has $\grad f=(0,0,0)$ at the origin, and hence the Pre-image Theorem fails there.
    \item Attempt to parametrise the cone near $(0,0,0)$ and show the regularity condition breaks down.
    \item Contrast with the hyperboloid $z^2=x^2+y^2+1$, where the Pre-image Theorem applies everywhere. Compare the two geometries.
  \end{enumerate}
\end{newq}

\begin{newq}
  \NEW\; \textbf{Surfaces as Manifolds (Section~2.7).}\;
  \begin{enumerate}[label=(\alph*)]
    \item State the definition of a smooth 2-dimensional manifold (as in Section~2.7), including the Hausdorff and second-countable conditions.
    \item Show that a regular surface $S\subset\R^3$ is a smooth manifold of dimension 2 by exhibiting a compatible atlas.
    \item Give an example of a smooth 2-manifold that cannot be embedded in $\R^3$ (e.g.\ the Klein bottle). Explain why regular surfaces in $\R^3$ are more restrictive.
  \end{enumerate}
\end{newq}

\begin{newq}
  \NEW\; \textbf{Coordinate Function Counter-Example.}\;
  The lecture notes give a counter-example: $\mathbf{x}(u,v)=(\sin 2u, \sin u, v)$ on $(-\pi,\pi)\times(-1,1)$ is a smooth bijection with $\rank(d\mathbf{x})=2$, yet $\mathbf{x}^{-1}$ is not continuous.
  \begin{enumerate}[label=(\alph*)]
    \item Describe the image surface $\mathbf{x}(U)$.
    \item Find an explicit sequence $\mathbf{x}(u_n,v_n)\to\mathbf{x}(0,0)$ but $(u_n,v_n)\not\to(0,0)$, demonstrating discontinuity of $\mathbf{x}^{-1}$.
    \item Explain why this does not contradict the regularity theorem (Proposition~4 in [DoCarmo]): what additional assumption is missing?
  \end{enumerate}
\end{newq}

\begin{newq}
  \NEW\; \textbf{Composition of Smooth Maps Between Surfaces.}\;
  The lecture notes prove: if $\varphi:S_1\to S_2$ and $\psi:S_2\to S_3$ are smooth, so is $\psi\circ\varphi$.
  \begin{enumerate}[label=(\alph*)]
    \item Reproduce the proof using coordinate functions.
    \item Identify where the chain rule is used in the argument.
    \item Give an explicit example: let $S_1=S_2=S_3=S^2$, $\varphi$ = antipodal map, $\psi$ = rotation by $\pi/2$ about the $z$-axis. Compute $(\psi\circ\varphi)(x,y,z)$.
  \end{enumerate}
\end{newq}

\begin{newq}
  \NEW\; \textbf{Graph Surface: Normal Vector and Area.}\;
  Let $S$ be the graph of $z=f(x,y)=x^2+y^2$ (paraboloid).
  \begin{enumerate}[label=(\alph*)]
    \item Using $\mathbf{x}(u,v)=(u,v,u^2+v^2)$, compute $E, F, G$.
    \item Find the unit normal $\mathbf{N}$ and verify it is indeed a unit vector.
    \item Set up the area integral for the region $\{x^2+y^2\leq R^2\}$ on the paraboloid and evaluate it.
  \end{enumerate}
\end{newq}

\begin{newq}
  \NEW\; \textbf{Creative: Restricting Smooth Functions.}\;
  The lecture notes state that the restriction of a smooth function $F:\R^3\to\R$ to a regular surface $S$ is smooth on $S$.
  \begin{enumerate}[label=(\alph*)]
    \item Let $S$ be the surface of revolution of $x=z^2+1$ about the $z$-axis. Show $S$ is a regular surface (use the Pre-image Theorem).
    \item Consider $f(x,y,z)=\dfrac{1}{x^2+2y^2}$. Determine an open set $V\supset S$ on which $f$ is well-defined and smooth; hence conclude $f|_S$ is smooth.
    \item Give an example of a smooth function on $S$ that is \emph{not} the restriction of any polynomial on $\R^3$.
  \end{enumerate}
\end{newq}

\begin{newq}
  \NEW\; \textbf{Creative: Latitude Circles and Parametrisation.}\;
  On the sphere $S^2$, a \emph{latitude circle} is $\{\theta=\theta_0\}$ in spherical coordinates $(\theta,\varphi)$ where $\theta\in(0,\pi)$.
  \begin{enumerate}[label=(\alph*)]
    \item Parametrise $S^2$ by $\mathbf{x}(\theta,\varphi)=(\sin\theta\cos\varphi,\sin\theta\sin\varphi,\cos\theta)$ and compute $E,F,G$.
    \item Show the parametrisation is orthogonal ($F=0$). What does this say about latitude and longitude circles?
    \item Compute the arc length of the latitude circle $\theta=\theta_0$ as $\varphi$ varies from $0$ to $2\pi$, and recover the familiar formula $2\pi\sin\theta_0$.
  \end{enumerate}
\end{newq}

%===========================================================
\section{Second Fundamental Form and Curvatures}
%===========================================================

\begin{question}
Show that the Weingarten map $dN_p: T_pS\to T_pS$ is self-adjoint with respect to the first fundamental form.
\end{question}

\begin{question}
Prove $K = \det(\mathrm{II})/\det(\mathrm{I})$.
\end{question}

\begin{question}
Compute $K$ and $H$ for the sphere $x^2+y^2+z^2=R^2$ using both the Weingarten map and the formula.
\end{question}

\begin{question}
Find the principal curvatures of the helicoid and show $H=0$.
\end{question}

\begin{question}
Prove: if every point of a connected surface $S$ is umbilic, then $S$ is a plane or a sphere.
\end{question}

\begin{question}
Compute $K$ for the Enneper surface and show $K\leq 0$ everywhere.
\end{question}

\begin{question}
Define elliptic, parabolic, and hyperbolic points. Show $K<0$ at a hyperbolic point with a geometric example.
\end{question}

\begin{question}
\textbf{Euler's Theorem.}\;
Derive $\kp_n(\theta) = \kp_1\cos^2\theta + \kp_2\sin^2\theta$.
\end{question}

\begin{question}
Let $S$ be the graph $z=f(x,y)$. Compute $K$ at a critical point and relate it to $\det[f_{ij}]$.
\end{question}

\begin{question}
\textbf{Meusnier's Theorem.}\;
Prove that all curves through $p$ with the same tangent direction have the same normal curvature.
\end{question}

\begin{question}
\textbf{Second Fundamental Form of the Torus.}\;
Compute $e,f,g$, $K$ and $H$ as functions of $u$. Identify where $K=0$.
\end{question}

\begin{question}
\textbf{Lines of Curvature.}\;
Show $\alpha$ is a line of curvature iff $dN(\alpha')$ is parallel to $\alpha'$.
\end{question}

\begin{question}
\textbf{Rodrigues' Formula.}\;
Prove a unit-speed curve is a line of curvature iff $N'(s)=-\kp_i\,\alpha'(s)$.
\end{question}

\begin{question}
\textbf{Asymptotic Curves.}\;
Show they exist only at non-elliptic points. Compute asymptotic directions on $z=xy$.
\end{question}

\begin{question}
\textbf{Minimal Surfaces.}\;
Show $H=0$ iff $\kp_1=-\kp_2$. Verify the catenoid is minimal.
\end{question}

\begin{question}
\textbf{Shape Operator and Curvature Lines.}\;
With $\kp_1\neq\kp_2$ everywhere, show principal directions define two orthogonal foliations.
\end{question}

\begin{question}
\textbf{Curvature of a Surface of Revolution.}\;
Derive formulas for $K$ and $H$ in terms of $f,g$ and their derivatives.
\end{question}

\begin{hardq}
  \textbf{Bonnet's Theorem.} \hard\;
  State Bonnet's Theorem and explain why both the Gauss and Codazzi--Mainardi equations are needed.
\end{hardq}

\begin{question}
\textbf{Gauss Map.}\;
Define $N:S\to S^2$. Show $dN_p=-\mathcal{W}_p$ and the Jacobian of $N$ equals $K$.
\end{question}

\begin{hardq}
  \textbf{Codazzi--Mainardi Equations.} \hard\;
  Derive them from $N_{uv}=N_{vu}$.
\end{hardq}

%---------- NEW QUESTIONS 21-40 ----------

\begin{newq}
  \NEW\; \textbf{Gauss Map: Smoothness.}\;
  The lecture notes define $N(\mathbf{x}(u,v))=\dfrac{\mathbf{x}_u\wedge\mathbf{x}_v}{|\mathbf{x}_u\wedge\mathbf{x}_v|}$.
  \begin{enumerate}[label=(\alph*)]
    \item Show that $N(\mathbf{x}(u,v))$ is a smooth function of $(u,v)$.
    \item Verify that $N\in S^2$, i.e.\ $|N|=1$.
    \item Show that $dN_p$ maps $T_pS$ to $T_{N(p)}S^2 = T_pS$ (identify these tangent spaces).
  \end{enumerate}
\end{newq}

\begin{newq}
  \NEW\; \textbf{Gauss Map of the Graph.}\;
  For the graph $z=f(x,y)$ with coordinate function $\mathbf{x}(u,v)=(u,v,f(u,v))$:
  \begin{enumerate}[label=(\alph*)]
    \item Compute $\mathbf{x}_u\wedge\mathbf{x}_v$ and hence the Gauss map $N(u,v)=\dfrac{(-f_u,-f_v,1)}{\sqrt{1+f_u^2+f_v^2}}$.
    \item Compute $N_u=dN(\mathbf{x}_u)$ and $N_v=dN(\mathbf{x}_v)$.
    \item At a critical point $f_u=f_v=0$, simplify $N$ and compute $e,f,g$ directly.
  \end{enumerate}
\end{newq}

\begin{newq}
  \NEW\; \textbf{Second Fundamental Form Coefficients.}\;
  The lecture notes define $e=\mathrm{II}(\mathbf{x}_u,\mathbf{x}_u)$, $f=\mathrm{II}(\mathbf{x}_u,\mathbf{x}_v)$, $g=\mathrm{II}(\mathbf{x}_v,\mathbf{x}_v)$.
  \begin{enumerate}[label=(\alph*)]
    \item Show that $e = -\langle N_u, \mathbf{x}_u\rangle = \langle N, \mathbf{x}_{uu}\rangle$.
    \item Use this to compute $e, f, g$ for the sphere $\mathbf{x}(u,v)=(R\sin v\cos u, R\sin v\sin u, R\cos v)$.
    \item Verify that $K = (eg-f^2)/(EG-F^2) = 1/R^2$.
  \end{enumerate}
\end{newq}

\begin{newq}
  \NEW\; \textbf{The Matrix of the Weingarten Map.}\;
  The lecture notes derive $[dN]_{BB} = -(EG-F^2)^{-1}\begin{pmatrix}E&F\\F&G\end{pmatrix}^{-1}\begin{pmatrix}e&f\\f&g\end{pmatrix}$.
  \begin{enumerate}[label=(\alph*)]
    \item Derive this matrix expression by expressing $dN(\mathbf{x}_u)$ and $dN(\mathbf{x}_v)$ in the basis $\{\mathbf{x}_u,\mathbf{x}_v\}$.
    \item Use it to express $K=\det(dN)$ and $H=-\tfrac{1}{2}\mathrm{tr}(dN)$ in terms of $e,f,g,E,F,G$.
    \item Verify the formulas for the unit sphere ($E=G=1$, $F=0$, $e=g=1/R$, $f=0$) giving $K=1/R^2$ and $H=1/R$.
  \end{enumerate}
\end{newq}

\begin{newq}
  \NEW\; \textbf{Self-Adjointness of $dN$.}\;
  \begin{enumerate}[label=(\alph*)]
    \item Prove $\langle dN(\mathbf{x}_u), \mathbf{x}_v\rangle = \langle \mathbf{x}_u, dN(\mathbf{x}_v)\rangle$ by differentiating $\langle N, \mathbf{x}_u\rangle = 0$ with respect to $v$ and $\langle N, \mathbf{x}_v\rangle=0$ with respect to $u$.
    \item Show this implies $f=\langle N_u, -\mathbf{x}_v\rangle = \langle N_v, -\mathbf{x}_u\rangle$, and hence the matrix $\begin{pmatrix}e&f\\f&g\end{pmatrix}$ is symmetric.
    \item Explain why self-adjointness implies $dN_p$ is diagonalisable with real eigenvalues.
  \end{enumerate}
\end{newq}

\begin{newq}
  \NEW\; \textbf{Normal Curvature from the Second Fundamental Form.}\;
  Let $\alpha(s)$ be a unit-speed curve on $S$ with $\alpha'(s)=\mathbf{v}$.
  \begin{enumerate}[label=(\alph*)]
    \item Show $\kp_n = \mathrm{II}(\mathbf{v},\mathbf{v}) = -\langle dN(\mathbf{v}),\mathbf{v}\rangle$.
    \item For a curve $\mathbf{v}=\cos\theta\,\mathbf{e}_1+\sin\theta\,\mathbf{e}_2$ where $\mathbf{e}_1,\mathbf{e}_2$ are principal directions, derive $\kp_n=\kp_1\cos^2\theta+\kp_2\sin^2\theta$.
    \item For the sphere of radius $R$, confirm $\kp_n=1/R$ for \emph{every} unit tangent direction.
  \end{enumerate}
\end{newq}

\begin{newq}
  \NEW\; \textbf{Diagonalising the Gauss Map.}\;
  The lecture notes (Section~3.6) show that $dN_p$ can be diagonalised.
  \begin{enumerate}[label=(\alph*)]
    \item Explain why the self-adjointness of $dN_p$ guarantees it is diagonalisable over $\R$.
    \item Let $\mathbf{e}_1,\mathbf{e}_2$ be orthonormal eigenvectors with eigenvalues $-\kp_1,-\kp_2$. Show $K=\kp_1\kp_2$ and $H=(\kp_1+\kp_2)/2$.
    \item For the Enneper surface at the origin, show directly that $dN_0$ has equal and opposite eigenvalues, confirming $H=0$.
  \end{enumerate}
\end{newq}

\begin{newq}
  \NEW\; \textbf{Gauss Map in Local Coordinates.}\;
  Section~3.7 derives the formula for $[dN]$ in local coordinates.
  \begin{enumerate}[label=(\alph*)]
    \item Verify: setting $a=1, b=0$ gives $dN(\mathbf{x}_u)=e\cdot\text{(first row of }(FFF)^{-1}\mathrm{II})$.
    \item For the torus $\mathbf{x}(u,v)=((R+r\cos u)\cos v,(R+r\cos u)\sin v, r\sin u)$, compute $N(u,v)$.
    \item Compute $N_u$ and $N_v$ directly and express them in the basis $\{\mathbf{x}_u,\mathbf{x}_v\}$ to find $e,f,g$.
  \end{enumerate}
\end{newq}

\begin{newq}
  \NEW\; \textbf{Umbilic Points and Totally Umbilic Surfaces.}\;
  A point $p$ is umbilic if $\kp_1(p)=\kp_2(p)$, equivalently $dN_p=\lambda\mathrm{Id}$.
  \begin{enumerate}[label=(\alph*)]
    \item Show that on the sphere every point is umbilic, with $\kp_1=\kp_2=1/R$.
    \item On a plane, every point is umbilic with $\kp_1=\kp_2=0$.
    \item Prove the converse: if $dN_p=\lambda(p)\mathrm{Id}$ for all $p\in S$ (connected), then $\lambda$ is constant and $S$ is a sphere or a plane. (\emph{Hint:} differentiate $N_u=\lambda\mathbf{x}_u$ and use Codazzi.)
  \end{enumerate}
\end{newq}

\begin{newq}
  \NEW\; \textbf{K at a Graph Critical Point.}\;
  Let $S$ be the graph of $z=f(x,y)$.  Suppose $f_x(0,0)=f_y(0,0)=0$.
  \begin{enumerate}[label=(\alph*)]
    \item Show $N(0,0)=(0,0,1)$ and compute the second fundamental form at $(0,0,0)$.
    \item Show $K(0,0,0) = f_{xx}f_{yy}-f_{xy}^2 = \det(\mathrm{Hess}_f(0,0))$.
    \item Apply this to $z=x^2-y^2$ (saddle) and $z=x^2+y^2$ (bowl) and classify the type of point.
  \end{enumerate}
\end{newq}

\begin{newq}
  \NEW\; \textbf{Linear Algebra Review for the Shape Operator.}\;
  Let $\mathcal{W}=-dN: T_pS\to T_pS$ be the shape operator (Weingarten map).
  \begin{enumerate}[label=(\alph*)]
    \item Show $\mathrm{II}(\mathbf{v},\mathbf{w})=\langle\mathcal{W}(\mathbf{v}),\mathbf{w}\rangle$ for all $\mathbf{v},\mathbf{w}\in T_pS$.
    \item Verify this symmetry using the formula $\mathrm{II}(\mathbf{v},\mathbf{w})=\langle N, \mathbf{v}_u w^1+\mathbf{v}_v w^2\rangle$ where appropriate.
    \item Explain why $\mathcal{W}$ being self-adjoint implies $\mathrm{II}$ is a symmetric bilinear form.
  \end{enumerate}
\end{newq}

\begin{newq}
  \NEW\; \textbf{Surface of Revolution: Curvatures.}\;
  For the surface $\mathbf{x}(u,v)=(f(v)\cos u, f(v)\sin u, g(v))$ with $(f')^2+(g')^2=1$ (unit-speed profile):
  \begin{enumerate}[label=(\alph*)]
    \item Compute $E,F,G$ and the unit normal $N$.
    \item Show that the principal directions at each point are along $\mathbf{x}_u$ and $\mathbf{x}_v$.
    \item Derive $\kp_1 = -g'/f$ (azimuthal) and $\kp_2 = f''g'-f'g''$ (meridional), and express $K$ and $H$.
  \end{enumerate}
\end{newq}

\begin{newq}
  \NEW\; \textbf{Catenoid is Minimal.}\;
  \begin{enumerate}[label=(\alph*)]
    \item For the catenoid $\mathbf{x}(u,v)=(c\cosh v\cos u, c\cosh v\sin u, cv)$, compute $e,f,g$.
    \item Verify $H=(eG-2fF+gE)/(2(EG-F^2))=0$.
    \item Show directly that $\kp_1=-\kp_2$ at every point.
  \end{enumerate}
\end{newq}

\begin{newq}
  \NEW\; \textbf{Asymptotic Directions on the Saddle.}\;
  \begin{enumerate}[label=(\alph*)]
    \item For $z=xy$ parametrised by $\mathbf{x}(u,v)=(u,v,uv)$, compute $e,f,g$.
    \item Find the asymptotic directions (solutions of $e(\dot u)^2+2f\dot u\dot v+g(\dot v)^2=0$).
    \item Verify the asymptotic curves are the lines $y=\text{const}$ and $x=\text{const}$ on the saddle surface.
  \end{enumerate}
\end{newq}

\begin{newq}
  \NEW\; \textbf{Creative: Soap Film and Mean Curvature.}\;
  A soap film minimises area subject to boundary conditions and satisfies $H=0$.
  \begin{enumerate}[label=(\alph*)]
    \item Explain physically why $H=0$ corresponds to mechanical equilibrium (equal pressure from both sides of the film).
    \item Show that for a surface of revolution $H=0$ reduces to the ODE $f''/(1+f'^2)^{3/2} = f'/(f\sqrt{1+f'^2})$, and verify the catenoid $f(z)=c\cosh(z/c)$ satisfies it.
    \item If you stretched a soap film between two coaxial circles of radius $r$ at heights $z=\pm h$, what happens as $h/r$ exceeds a certain ratio? (Describe qualitatively; relate to $H=0$ losing solutions.)
  \end{enumerate}
\end{newq}

\begin{newq}
  \NEW\; \textbf{Creative: Curvature of the Earth.}\;
  Model the Earth as a sphere of radius $R=6400$ km.
  \begin{enumerate}[label=(\alph*)]
    \item What is the Gaussian curvature $K$ of the Earth's surface?
    \item If you travel along the equator (a great circle), what is the normal curvature $\kp_n$ and geodesic curvature $\kp_g$?
    \item A flat map of a region of area $A$ on the Earth's surface will have distortion. Using the Theorema Egregium, explain why no flat map can be perfect.
  \end{enumerate}
\end{newq}

\begin{newq}
  \NEW\; \textbf{Codazzi--Mainardi: Physical Interpretation.}\;
  The Codazzi--Mainardi equations are integrability conditions on the second fundamental form.
  \begin{enumerate}[label=(\alph*)]
    \item Write out the Codazzi--Mainardi equations in the orthogonal case $F=0$.
    \item Explain: they ensure that $N_u$ and $N_v$ computed from $dN$ are consistent as mixed partials ($N_{uv}=N_{vu}$).
    \item Give an example: verify Codazzi--Mainardi for the sphere, where $e=g=1/R$, $f=0$.
  \end{enumerate}
\end{newq}

\begin{newq}
  \NEW\; \textbf{Creative: Sign of Gaussian Curvature.}\;
  Consider the torus $\mathbf{x}(u,v)=((R+r\cos u)\cos v,(R+r\cos u)\sin v,r\sin u)$ with $R>r>0$.
  \begin{enumerate}[label=(\alph*)]
    \item Show $K(u,v) = \dfrac{\cos u}{r(R+r\cos u)}$.
    \item Identify the outer equator ($u=0$), inner equator ($u=\pi$), and top/bottom circles ($u=\pi/2, 3\pi/2$) as elliptic, hyperbolic, or parabolic points.
    \item Using Gauss--Bonnet, explain why $\iint_T K\,dA=0$ is consistent with your sign analysis.
  \end{enumerate}
\end{newq}

\begin{newq}
  \NEW\; \textbf{Creative: Ruled Surfaces and Asymptotic Lines.}\;
  A \emph{ruled surface} is swept out by a moving line: $\mathbf{x}(u,v)=\mathbf{a}(u)+v\mathbf{b}(u)$.
  \begin{enumerate}[label=(\alph*)]
    \item Show that along each ruling $v\mapsto\mathbf{a}(u_0)+v\mathbf{b}(u_0)$, the normal curvature in the direction of $\mathbf{b}$ is zero.
    \item Conclude that the rulings are asymptotic curves, so $K\leq 0$ at all non-parabolic points.
    \item The helicoid is a ruled surface. Verify $K\leq 0$ explicitly from your earlier computation of $K$ for the helicoid.
  \end{enumerate}
\end{newq}

\begin{newq}
  \NEW\; \textbf{Gauss Map Area Ratio.}\;
  The Gauss map $N:S\to S^2$ carries area elements in the ratio $|K|$.
  \begin{enumerate}[label=(\alph*)]
    \item Show: the area element pulled back by $N$ satisfies $N^*dA_{S^2} = K\,dA_S$.
    \item For the sphere $S^2_R$, verify that the total image area under $N$ is $4\pi$ (the full unit sphere), consistent with $K=1/R^2$ and $\text{Area}(S^2_R)=4\pi R^2$.
    \item For a cylinder $x^2+y^2=r^2$, explain why the Gauss map sends all points to a single great circle, and reconcile with $K=0$.
  \end{enumerate}
\end{newq}

%===========================================================
\section{Isometries and Theorema Egregium}
%===========================================================

\begin{question}
Define a local isometry. Show $E_1=E_2, F_1=F_2, G_1=G_2$ is sufficient for two surfaces to be locally isometric.
\end{question}

\begin{question}
Show the sphere and plane are not locally isometric by comparing $K$, and explain the consequence for flat maps.
\end{question}

\begin{question}
Show the catenoid and helicoid are locally isometric via an explicit family of isometries.
\end{question}

\begin{question}
\textbf{Theorema Egregium.}\;
State the theorem and explain why $K$ depends only on $E,F,G$.
\end{question}

\begin{question}
\textbf{Developable Surfaces.}\;
Prove $K\equiv 0$ implies locally isometric to a plane. Give three classes of examples.
\end{question}

\begin{question}
Show if $\lambda\equiv 1$ in a conformal map, then $f$ is a local isometry.
\end{question}

\begin{question}
Show the cylinder $x^2+y^2=r^2$ is locally isometric to a plane.
\end{question}

\begin{question}
Compute Christoffel symbols $\Gamma_{ij}^k$ for $S^2$ in spherical coordinates.
\end{question}

\begin{hardq}
  \textbf{Brioschi Formula.} \hard\;
  Derive the Brioschi formula for $K$ when $F=0$ and verify for the sphere and cylinder.
\end{hardq}

\begin{question}
Give a concrete example of an equiareal map that is not an isometry. What additional condition forces isometry?
\end{question}

\begin{question}
\textbf{Rigidity of the Sphere.}\;
State the rigidity theorem and contrast with non-convex surfaces.
\end{question}

\begin{question}
\textbf{Gauss Equations.}\;
Write out $R_{ijk\ell}$ in terms of the second fundamental form and recover $K=(eg-f^2)/(EG-F^2)$.
\end{question}

\begin{question}
\textbf{Isometry Group.}\;
Show isometries $S\to S$ form a group and identify the isometry group of $S^2$.
\end{question}

\begin{question}
\textbf{Conformal Flatness.}\;
State (without full proof) that every regular surface admits isothermal coordinates.
\end{question}

\begin{question}
\textbf{Intrinsic vs.\ Extrinsic.}\;
Classify: (a) $K$, (b) $H$, (c) arc length, (d) $\kp_1,\kp_2$, (e) Christoffel symbols.
\end{question}

\begin{question}
Compute $K$ for the pseudosphere and verify $K=-1$ via the Brioschi formula.
\end{question}

\begin{question}
\textbf{Nash Embedding Theorem.}\;
State the theorem and explain its implication for abstract Riemannian metrics on 2-manifolds.
\end{question}

\begin{hardq}
  \textbf{Parallel Transport and Curvature.} \hard\;
  Show $K$ can be recovered as holonomy angle per unit area and relate to Gauss--Bonnet.
\end{hardq}

\begin{question}
\textbf{Covariant Derivative.}\;
Define $\nabla_X Y$ on a surface. State the properties of the Levi-Civita connection.
\end{question}

\begin{hardq}
  \textbf{Minding's Theorem.} \hard\;
  State and use it to conclude all constant negative curvature surfaces share the same local geometry.
\end{hardq}

%---------- NEW QUESTIONS 21-40 ----------

\begin{newq}
  \NEW\; \textbf{Local Isometry and the First Fundamental Form.}\;
  \begin{enumerate}[label=(\alph*)]
    \item Define a local isometry $\varphi: S_1\to S_2$ precisely.
    \item Show: if $\varphi$ is a local isometry, then $E_1=E_2\circ\varphi$, $F_1=F_2\circ\varphi$, $G_1=G_2\circ\varphi$ in appropriate parametrisations.
    \item Show the converse: if $E_1=E_2$, $F_1=F_2$, $G_1=G_2$, then the map induced by the coordinate functions is a local isometry.
  \end{enumerate}
\end{newq}

\begin{newq}
  \NEW\; \textbf{Rigid Motions of $\R^3$ Are Isometries.}\;
  Let $R(\mathbf{x})=A\mathbf{x}+\mathbf{v}$ be a rigid motion with $A\in SO(3)$.
  \begin{enumerate}[label=(\alph*)]
    \item Show $R$ induces an isometry between any surface $S$ and its image $R(S)$.
    \item Show that rigid motions preserve the first fundamental form.
    \item Conclude they also preserve $K$, $H$, $\kp_1$, $\kp_2$.
  \end{enumerate}
\end{newq}

\begin{newq}
  \NEW\; \textbf{Gaussian Curvature Is Intrinsic.}\;
  \begin{enumerate}[label=(\alph*)]
    \item State the Theorema Egregium precisely.
    \item The Brioschi formula expresses $K$ in terms of $E,F,G$ and their partial derivatives only. For a surface with $F=0$, it reads
    \[K = -\frac{1}{2\sqrt{EG}}\left[\frac{\partial}{\partial u}\!\left(\frac{G_u}{\sqrt{EG}}\right)+\frac{\partial}{\partial v}\!\left(\frac{E_v}{\sqrt{EG}}\right)\right].\]
    Verify this for the unit sphere (spherical coordinates, $E=1, F=0, G=\sin^2\theta$) and confirm $K=1$.
    \item Explain in one sentence why $H$ is \emph{not} intrinsic.
  \end{enumerate}
\end{newq}

\begin{newq}
  \NEW\; \textbf{Cylinder Is Locally Isometric to a Plane.}\;
  \begin{enumerate}[label=(\alph*)]
    \item Write the coordinate function for the cylinder and the plane and show their first fundamental forms coincide.
    \item Exhibit the explicit local isometry $\varphi: \text{cylinder}\to\text{plane}$.
    \item Verify $K_{\text{cylinder}}=0$ two ways: (i) from the second fundamental form; (ii) from the Theorema Egregium (since the plane has $K=0$ and the map is an isometry).
  \end{enumerate}
\end{newq}

\begin{newq}
  \NEW\; \textbf{Christoffel Symbols from the First Fundamental Form.}\;
  \begin{enumerate}[label=(\alph*)]
    \item Derive the general formula $\Gamma_{ij}^k = \tfrac{1}{2}g^{kl}(\partial_i g_{jl}+\partial_j g_{il}-\partial_l g_{ij})$ where $g_{ij}=\begin{pmatrix}E&F\\F&G\end{pmatrix}$.
    \item Compute all $\Gamma_{ij}^k$ for the sphere in spherical coordinates $(\theta,\varphi)$ where $E=1, F=0, G=\sin^2\theta$.
    \item Verify: $\Gamma_{22}^1=-\sin\theta\cos\theta$ and $\Gamma_{12}^2=\cot\theta$, all others being zero or equivalent.
  \end{enumerate}
\end{newq}

\begin{newq}
  \NEW\; \textbf{Intrinsic vs.\ Extrinsic Quantities.}\;
  \begin{enumerate}[label=(\alph*)]
    \item For each quantity: $K$, $H$, arc length, $\kp_1$, $\kp_2$, Christoffel symbols $\Gamma_{ij}^k$ --- state whether it is intrinsic or extrinsic and justify briefly.
    \item Show that $K=\kp_1\kp_2$ is intrinsic even though $\kp_1$ and $\kp_2$ individually are extrinsic. (\emph{Hint:} $K$ depends only on $E,F,G$.)
    \item Give an example of two surfaces that are locally isometric but have different mean curvatures everywhere.
  \end{enumerate}
\end{newq}

\begin{newq}
  \NEW\; \textbf{Conformal Map vs.\ Isometry.}\;
  \begin{enumerate}[label=(\alph*)]
    \item Define a conformal map. Show that every isometry is conformal (with $\lambda=1$).
    \item Give an example of a conformal map that is not an isometry: stereographic projection $S^2\setminus\{N\}\to\R^2$.
    \item Compute the conformal factor $\lambda$ for stereographic projection and verify $\lambda\neq 1$ in general.
  \end{enumerate}
\end{newq}

\begin{newq}
  \NEW\; \textbf{Developable Surfaces.}\;
  \begin{enumerate}[label=(\alph*)]
    \item Prove: a surface with $K\equiv 0$ is locally isometric to a plane (use the Theorema Egregium and the fact that flat metrics are locally Euclidean).
    \item List three classes of developable surfaces (plane, cylinder, cone) and show each has $K=0$.
    \item Explain why a cone is locally isometric to a plane but not globally, by describing what happens to the ``apex'' under the development.
  \end{enumerate}
\end{newq}

\begin{newq}
  \NEW\; \textbf{Equiareal Maps vs.\ Isometries.}\;
  \begin{enumerate}[label=(\alph*)]
    \item Show that an isometry preserves area, hence is equiareal.
    \item Construct an explicit equiareal map from the sphere to a cylinder that is not an isometry. (\emph{Hint:} Lambert's equal-area projection.)
    \item Identify which quantity (preserved vs.\ changed) distinguishes the two maps.
  \end{enumerate}
\end{newq}

\begin{newq}
  \NEW\; \textbf{The Gauss Equation.}\;
  The Gauss equation relates the Riemann curvature to the second fundamental form:
  \[R^1{}_{212} = eg - f^2.\]
  \begin{enumerate}[label=(\alph*)]
    \item State the Gauss equation in the form $R_{1212}=eg-f^2$ and explain what $R_{1212}$ measures.
    \item Use it to recover $K=(eg-f^2)/(EG-F^2)$ and explain why this proves $K$ is intrinsic.
    \item Verify for the sphere: compute $R_{1212}$ from Christoffel symbols and confirm $R_{1212}=\sin^2\theta = eg-f^2$.
  \end{enumerate}
\end{newq}

\begin{newq}
  \NEW\; \textbf{Covariant Derivative on a Surface.}\;
  The covariant derivative $DV/dt$ of a vector field $V$ along a curve $\gamma$ is the projection of $V'$ onto $T_{\gamma(t)}S$.
  \begin{enumerate}[label=(\alph*)]
    \item Show: $\dfrac{DV}{dt}=\dfrac{dV}{dt}-\left\langle\dfrac{dV}{dt},N\right\rangle N$.
    \item Prove: $\dfrac{d}{dt}\langle V,W\rangle = \left\langle\dfrac{DV}{dt},W\right\rangle+\left\langle V,\dfrac{DW}{dt}\right\rangle$ (metric compatibility).
    \item Show the Levi-Civita connection is torsion-free: $D_{\mathbf{x}_u}\mathbf{x}_v - D_{\mathbf{x}_v}\mathbf{x}_u = 0$.
  \end{enumerate}
\end{newq}

\begin{newq}
  \NEW\; \textbf{Pseudosphere and Constant Negative Curvature.}\;
  The pseudosphere has $\mathbf{x}(u,v)=(\mathrm{sech}\,v\cos u,\mathrm{sech}\,v\sin u, v-\tanh v)$.
  \begin{enumerate}[label=(\alph*)]
    \item Compute $E, F, G$ and verify $F=0$.
    \item Apply the Brioschi formula to confirm $K=-1$.
    \item State Minding's Theorem and use it to conclude the pseudosphere is locally isometric to the Poincar\'e half-plane (which also has $K=-1$).
  \end{enumerate}
\end{newq}

\begin{newq}
  \NEW\; \textbf{Holonomy as Curvature.}\;
  \begin{enumerate}[label=(\alph*)]
    \item Define holonomy: the angle $\Delta\theta$ by which a vector rotates when parallel-transported around a closed curve.
    \item For a geodesic triangle on $S^2$ with area $A$, show the holonomy is $A/R^2$ (in radians).
    \item For a small geodesic rectangle of area $\Delta A$ on a surface with $K=K_0$, show $\Delta\theta\approx K_0\,\Delta A$.
  \end{enumerate}
\end{newq}

\begin{newq}
  \NEW\; \textbf{Catenoid--Helicoid Isometry.}\;
  \begin{enumerate}[label=(\alph*)]
    \item Parametrise the helicoid $\mathbf{x}(u,v)=(v\cos u, v\sin u, cu)$ and the catenoid $\mathbf{y}(u,v)=(c\cosh v\cos u, c\cosh v\sin u, cv)$. Compute their first fundamental forms.
    \item Show they are equal, hence the natural map $\mathbf{y}\circ\mathbf{x}^{-1}$ is a local isometry.
    \item Since both $K_{\text{helicoid}}$ and $K_{\text{catenoid}}$ can be computed, verify they agree (by Theorema Egregium they must). Compute $K$ explicitly.
  \end{enumerate}
\end{newq}

\begin{newq}
  \NEW\; \textbf{Isometry Group of Surfaces.}\;
  \begin{enumerate}[label=(\alph*)]
    \item Show that the composition of two isometries is an isometry and the inverse of an isometry is an isometry.
    \item Identify the isometry group of the flat cylinder $\{x^2+y^2=1\}$: describe all isometries explicitly.
    \item Show the isometry group of $S^2$ is $O(3)$. Why is it not just $SO(3)$?
  \end{enumerate}
\end{newq}

\begin{newq}
  \NEW\; \textbf{Creative: Paper Folding and Isometries.}\;
  When you fold a piece of paper, you perform an isometry (no stretching or tearing).
  \begin{enumerate}[label=(\alph*)]
    \item Explain why folding preserves $E,F,G$ and hence arc lengths and angles on the paper.
    \item By Theorema Egregium, $K$ is preserved. Since flat paper has $K=0$, what can you conclude about all surfaces obtained by folding?
    \item Can you fold a flat sheet into a doubly curved surface (positive Gaussian curvature, like a sphere cap)? Justify rigorously.
  \end{enumerate}
\end{newq}

\begin{newq}
  \NEW\; \textbf{Creative: Map Projections.}\;
  No flat map of the Earth can perfectly preserve distances.
  \begin{enumerate}[label=(\alph*)]
    \item Use the Theorema Egregium to give a rigorous proof: $K_{\text{sphere}}\neq K_{\text{plane}}$ implies no local isometry exists.
    \item The Mercator projection is conformal. What does it preserve and what does it distort?
    \item The equal-area projection preserves area. Write down the defining property as a condition on the first fundamental forms.
  \end{enumerate}
\end{newq}

\begin{newq}
  \NEW\; \textbf{Levi-Civita Connection: Uniqueness.}\;
  The Levi-Civita connection is the unique connection that is metric-compatible and torsion-free.
  \begin{enumerate}[label=(\alph*)]
    \item State both conditions precisely.
    \item Derive the Koszul formula:
    $2\langle\nabla_X Y, Z\rangle = X\langle Y,Z\rangle + Y\langle X,Z\rangle - Z\langle X,Y\rangle + \langle[X,Y],Z\rangle - \langle[Y,Z],X\rangle + \langle[X,Z],Y\rangle$.
    \item Show that the Christoffel symbols are the components of $\nabla_{\mathbf{x}_i}\mathbf{x}_j$ and use the Koszul formula to recover the formula $\Gamma_{ij}^k=\tfrac{1}{2}g^{kl}(\partial_i g_{jl}+\partial_j g_{il}-\partial_l g_{ij})$.
  \end{enumerate}
\end{newq}

\begin{newq}
  \NEW\; \textbf{Creative: Intrinsic Geometry of the Flat Torus.}\;
  The flat torus $T^2=\R^2/\mathbb{Z}^2$ has $E=G=1$, $F=0$ everywhere.
  \begin{enumerate}[label=(\alph*)]
    \item Show $K=0$ on the flat torus, hence it is locally isometric to a plane.
    \item Is the flat torus globally isometric to any open subset of $\R^2$? Justify topologically.
    \item The standard torus in $\R^3$ has $K$ varying (sometimes positive, sometimes negative). Explain why the standard torus in $\R^3$ is NOT isometric to the flat torus, even though both are tori. What is the obstruction?
  \end{enumerate}
\end{newq}

\begin{newq}
  \NEW\; \textbf{Creative: GPS and Curvature.}\;
  A GPS satellite measures distances along the curved Earth surface.
  \begin{enumerate}[label=(\alph*)]
    \item Model the Earth as a sphere of radius $R$. Compute the geodesic distance between two points $(0,0)$ and $(0,\Delta\varphi)$ in spherical coordinates.
    \item By Theorema Egregium, the GPS must account for the curvature of space. Estimate the error in flat-Earth distance (straight-line) vs.\ the true geodesic distance for $\Delta\varphi = 0.01$ radians.
    \item How does the answer change if the Earth were a prolate spheroid (football shape)? Which has larger $K$?
  \end{enumerate}
\end{newq}

%===========================================================
\section{Geodesics and Parallel Transport}
%===========================================================

\begin{question}
Derive the geodesic equations $\ddot{u}^k + \Gamma_{ij}^k\dot{u}^i\dot{u}^j=0$ using the calculus of variations.
\end{question}

\begin{question}
\textbf{Geodesics on the Cylinder.}\;
Show geodesics on $x^2+y^2=r^2$ are helices.
\end{question}

\begin{question}
Show great circles are geodesics on $S^2$ via the geodesic equations in spherical coordinates.
\end{question}

\begin{question}
\textbf{Clairaut's Relation.}\;
Prove $r\sin\psi=\text{const}$ for a geodesic on a surface of revolution.
\end{question}

\begin{question}
Define geodesic curvature $\kp_g$. Show $\gamma$ is geodesic iff $\kp_g\equiv 0$. Relate to $\kp^2=\kp_n^2+\kp_g^2$.
\end{question}

\begin{question}
Compute all geodesics on the Poincar\'e half-plane. Show they are semicircles centred on the $x$-axis and vertical lines.
\end{question}

\begin{question}
Define parallel transport. Show transporting a vector around a spherical triangle of area $\Omega$ rotates it by $\Omega$.
\end{question}

\begin{question}
\textbf{Gauss--Bonnet Theorem.}\;
State the global theorem and use it to find the angle sum in a geodesic triangle on $S^2$ with area $A$.
\end{question}

\begin{question}
Show $DV/dt$ is the projection of $V'$ onto $T_{\gamma(t)}S$.
\end{question}

\begin{hardq}
  \textbf{Exponential Map.} \hard\;
  Define $\exp_p$ and show it is a local diffeomorphism. Define normal coordinates and show $g_{ij}(p)=\delta_{ij}$, $\Gamma_{ij}^k(p)=0$.
\end{hardq}

\begin{question}
\textbf{Geodesics on the Flat Torus.}\;
Show geodesics on $\R^2/\mathbb{Z}^2$ are images of straight lines. Classify closed geodesics.
\end{question}

\begin{question}
\textbf{Geodesic Completeness.}\;
Define it. State Hopf--Rinow and deduce that every compact regular surface is geodesically complete.
\end{question}

\begin{question}
\textbf{Cut Locus.}\;
Define the cut locus $\mathrm{Cut}(p)$. Describe $\mathrm{Cut}(\text{north pole})$ on $S^2$.
\end{question}

\begin{question}
\textbf{Local Gauss--Bonnet.}\;
State the local formula and apply it to find the angle sum for a geodesic triangle with constant curvature $k$.
\end{question}

\begin{question}
\textbf{Jacobi Fields.}\;
Define a Jacobi field and show it arises as a variation through geodesics.
\end{question}

\begin{question}
\textbf{Conjugate Points.}\;
Define them. Show on $S^2_R$ every geodesic has conjugate points at distance $\pi R$.
\end{question}

\begin{hardq}
  \textbf{Bonnet's Diameter Theorem.} \hard\;
  State: $K\geq\delta>0$ complete $\Rightarrow$ $\mathrm{diam}(S)\leq\pi/\sqrt{\delta}$. Sketch the Jacobi field argument.
\end{hardq}

\begin{question}
\textbf{Index Form.}\;
Define $I(V,W)=\int_a^b(\langle V',W'\rangle - K\langle V,W\rangle)dt$. State when $\gamma$ minimises length locally.
\end{question}

\begin{question}
\textbf{Sphere Theorem.}\;
State: compact orientable surface with $K>0$ is homeomorphic to $S^2$. Explain why $K>0$ without compactness is insufficient.
\end{question}

\begin{hardq}
  \textbf{Euler Characteristic and Topology.} \hard\;
  Using Gauss--Bonnet, prove no metric on $T^2$ can have $K>0$ everywhere, and no metric on $S^2$ can have $K\leq 0$. Generalise to genus $g$.
\end{hardq}

%---------- NEW QUESTIONS 21-40 ----------

\begin{newq}
  \NEW\; \textbf{Geodesic Equations from Energy Minimisation.}\;
  A geodesic minimises the energy functional $E[\gamma]=\tfrac{1}{2}\int_a^b g(\gamma',\gamma')\,dt$.
  \begin{enumerate}[label=(\alph*)]
    \item Set up the Euler--Lagrange equations for $E[\gamma]$ and derive the geodesic equations $\ddot u^k+\Gamma_{ij}^k\dot u^i\dot u^j=0$.
    \item Explain why minimising energy (rather than arc length) produces the same geodesics but with constant speed automatically.
    \item Verify: the geodesic equations on the plane ($E=G=1$, $F=0$, all $\Gamma=0$) give straight lines $\ddot u^k=0$.
  \end{enumerate}
\end{newq}

\begin{newq}
  \NEW\; \textbf{Geodesics Have Constant Speed.}\;
  \begin{enumerate}[label=(\alph*)]
    \item Show $\dfrac{d}{dt}|\gamma'|^2 = 0$ along any geodesic, so speed is constant.
    \item Use this to show that arc-length parametrised geodesics ($|\gamma'|=1$) satisfy the standard geodesic equations.
    \item For the sphere, great circles traversed at constant angular speed have $|\gamma'|=R$; verify this is consistent with the geodesic equations.
  \end{enumerate}
\end{newq}

\begin{newq}
  \NEW\; \textbf{Great Circles on the Sphere.}\;
  In spherical coordinates $\mathbf{x}(\theta,\varphi)=(\sin\theta\cos\varphi,\sin\theta\sin\varphi,\cos\theta)$:
  \begin{enumerate}[label=(\alph*)]
    \item Write out the geodesic equations using the Christoffel symbols from Section~4.
    \item Verify that the equator $\theta(t)=\pi/2$, $\varphi(t)=t$ satisfies the geodesic equations.
    \item Show that any great circle can be obtained from the equator by a rigid rotation, hence every great circle is a geodesic.
  \end{enumerate}
\end{newq}

\begin{newq}
  \NEW\; \textbf{Clairaut's Relation on the Torus.}\;
  Apply Clairaut's relation $r\sin\psi=c$ to the torus with $r=R+r_0\cos u$.
  \begin{enumerate}[label=(\alph*)]
    \item Identify which latitude circles $u=u_0$ are geodesics. (\emph{Hint:} For a latitude circle, $\psi=\pi/2$.)
    \item For $c=R$ (outermost equator), describe the geodesic qualitatively. Does it wrap around the torus or stay near the equator?
    \item Show that geodesics starting from the innermost circle $r=R-r_0$ cannot penetrate into the inner region where $r<c$.
  \end{enumerate}
\end{newq}

\begin{newq}
  \NEW\; \textbf{Geodesic Curvature Formula.}\;
  For a unit-speed curve $\alpha(s)$ on a surface $S$ with unit normal $N$:
  \begin{enumerate}[label=(\alph*)]
    \item Show $\kp_g = \langle\alpha''\!, N\wedge\alpha'\rangle$.
    \item Verify $\kp^2=\kp_n^2+\kp_g^2$ using $\alpha''=\kp_g(N\wedge\alpha')+\kp_n N$.
    \item Compute $\kp_g$ for a latitude circle $\theta=\theta_0$ on the unit sphere and verify the formula gives $\kp_g=\cot\theta_0$.
  \end{enumerate}
\end{newq}

\begin{newq}
  \NEW\; \textbf{Parallel Transport Preserves Inner Products.}\;
  Let $V(t)$ and $W(t)$ be parallel vector fields along a curve $\gamma(t)$ on $S$.
  \begin{enumerate}[label=(\alph*)]
    \item Show $\dfrac{d}{dt}\langle V,W\rangle = 0$.
    \item Conclude that parallel transport is an orthogonal transformation $T_{\gamma(a)}S\to T_{\gamma(b)}S$.
    \item For the sphere, show that parallel transporting a vector along a latitude circle $\theta=\theta_0$ from $\varphi=0$ to $\varphi=2\pi$ rotates it by angle $2\pi\cos\theta_0$.
  \end{enumerate}
\end{newq}

\begin{newq}
  \NEW\; \textbf{Gauss--Bonnet for a Geodesic Triangle.}\;
  Let $T$ be a geodesic triangle on a surface $S$ with interior angles $\alpha_1,\alpha_2,\alpha_3$.
  \begin{enumerate}[label=(\alph*)]
    \item State the local Gauss--Bonnet formula for $T$: $\sum_{i=1}^3\alpha_i = \pi + \iint_T K\,dA$.
    \item For the unit sphere ($K=1$), a triangle with area $A$ has angle sum $\pi+A$. Verify for the octant triangle $(\alpha_1=\alpha_2=\alpha_3=\pi/2$, $A=\pi/2$).
    \item For a flat surface ($K=0$), confirm the angle sum is $\pi$ (Euclidean geometry).
  \end{enumerate}
\end{newq}

\begin{newq}
  \NEW\; \textbf{Covariant Derivative and Projection.}\;
  \begin{enumerate}[label=(\alph*)]
    \item Show that for a vector field $V$ along $\gamma$ on $S$, the covariant derivative is $DV/dt = V' - \langle V', N\rangle N$.
    \item A vector field is \emph{parallel} if $DV/dt=0$. Show this means $V'$ is always normal to $S$.
    \item For the sphere, show that the tangent vector $\gamma'$ of a great circle is parallel along itself ($D\gamma'/dt=0$), confirming great circles are geodesics.
  \end{enumerate}
\end{newq}

\begin{newq}
  \NEW\; \textbf{Exponential Map and Normal Coordinates.}\;
  \begin{enumerate}[label=(\alph*)]
    \item Define $\exp_p: T_pS\to S$ by $\exp_p(\mathbf{v})=\gamma(1)$ where $\gamma$ is the geodesic with $\gamma(0)=p$, $\gamma'(0)=\mathbf{v}$.
    \item Show $d(\exp_p)_{\mathbf{0}}=\mathrm{id}$, hence $\exp_p$ is a local diffeomorphism near $\mathbf{0}\in T_pS$.
    \item In normal coordinates $(u^1,u^2)$ defined by $\exp_p(u^i\mathbf{e}_i)$, verify $g_{ij}(p)=\delta_{ij}$ and $\Gamma_{ij}^k(p)=0$.
  \end{enumerate}
\end{newq}

\begin{newq}
  \NEW\; \textbf{Jacobi Equation.}\;
  Let $\gamma_s(t)$ be a one-parameter family of geodesics. The \emph{Jacobi field} is $J(t)=\partial\gamma_s/\partial s|_{s=0}$.
  \begin{enumerate}[label=(\alph*)]
    \item Derive the Jacobi equation $J''+KJ=0$ (for surfaces).
    \item For $S^2$ (K=1), solve to get $J(t)=A\sin t + B\cos t$.  If $J(0)=0$ and $J'(0)\neq 0$, show the Jacobi field vanishes again at $t=\pi$.
    \item Conclude that geodesics on $S^2$ have conjugate points at distance $\pi$ and explain why this means geodesics longer than $\pi R$ are not minimising.
  \end{enumerate}
\end{newq}

\begin{newq}
  \NEW\; \textbf{Geodesics on the Flat Torus.}\;
  The flat torus is $\R^2/\mathbb{Z}^2$ with the Euclidean metric.
  \begin{enumerate}[label=(\alph*)]
    \item Show the Christoffel symbols vanish everywhere, so the geodesic equations become $\ddot u^k=0$.
    \item Conclude all geodesics are images of straight lines $u^k(t)=a^k t+b^k$.
    \item A geodesic is \emph{closed} iff $a^1/a^2\in\mathbb{Q}$. Verify this and count the closed geodesics up to homotopy.
  \end{enumerate}
\end{newq}

\begin{newq}
  \NEW\; \textbf{Global Gauss--Bonnet.}\;
  \begin{enumerate}[label=(\alph*)]
    \item State the global Gauss--Bonnet theorem: $\iint_S K\,dA = 2\pi\chi(S)$.
    \item Compute $\iint_S K\,dA$ for the sphere ($\chi=2$), torus ($\chi=0$), and genus-$g$ surface ($\chi=2-2g$).
    \item Deduce: if $K>0$ everywhere on a compact surface, then $\chi(S)>0$, so $S$ is homeomorphic to $S^2$.
  \end{enumerate}
\end{newq}

\begin{newq}
  \NEW\; \textbf{Hopf--Rinow and Geodesic Completeness.}\;
  \begin{enumerate}[label=(\alph*)]
    \item Define geodesic completeness (every geodesic can be extended indefinitely).
    \item State the Hopf--Rinow theorem and list its equivalent formulations.
    \item Prove: every compact regular surface $S\subset\R^3$ is geodesically complete. (\emph{Hint:} A geodesic that cannot be extended must leave every compact set, but $S$ is closed and bounded in $\R^3$.)
  \end{enumerate}
\end{newq}

\begin{newq}
  \NEW\; \textbf{Local vs.\ Global Minimisation.}\;
  A geodesic is \emph{locally minimising} but not necessarily globally minimising.
  \begin{enumerate}[label=(\alph*)]
    \item Show that every geodesic is locally minimising (has no conjugate points on short segments).
    \item For $S^2$, show that a geodesic arc of length $>\pi R$ is not the shortest path between its endpoints.
    \item Relate the cut locus $\mathrm{Cut}(p)$ to the set of points beyond which a geodesic from $p$ fails to minimise.
  \end{enumerate}
\end{newq}

\begin{newq}
  \NEW\; \textbf{Geodesics on the Cylinder via Unrolling.}\;
  \begin{enumerate}[label=(\alph*)]
    \item Show that ``unrolling'' the cylinder to a flat strip is an isometry.
    \item Since geodesics are preserved by isometries and the plane's geodesics are straight lines, conclude the cylinder's geodesics are helices (when re-rolled).
    \item Parametrise all geodesics on the cylinder explicitly and identify the special cases: vertical lines, horizontal circles, and general helices.
  \end{enumerate}
\end{newq}

\begin{newq}
  \NEW\; \textbf{Creative: Ant on a Sphere.}\;
  An ant walks along a geodesic on the unit sphere, starting at the north pole.
  \begin{enumerate}[label=(\alph*)]
    \item Show that any path starting from the north pole in a fixed direction traces a meridian (longitude circle).
    \item After walking a distance $\pi/2$ (quarter circle), the ant arrives at the equator. How far has it walked (in terms of arc length)?
    \item If two ants start from the north pole along meridians $\varphi=0$ and $\varphi=\pi/2$, at what point do their paths intersect again? What is the area of the geodesic triangle they bound?
  \end{enumerate}
\end{newq}

\begin{newq}
  \NEW\; \textbf{Creative: Foucault Pendulum and Parallel Transport.}\;
  A Foucault pendulum at colatitude $\theta_0$ swings in a vertical plane that slowly rotates.
  \begin{enumerate}[label=(\alph*)]
    \item Model the Earth as a sphere. The pendulum's swing direction is a vector being parallel-transported along the latitude circle $\theta=\theta_0$ as the Earth rotates.
    \item Using the holonomy formula for a latitude circle, show the pendulum's plane rotates by $2\pi\cos\theta_0$ per day.
    \item At the North Pole ($\theta_0=0$), at the equator ($\theta_0=\pi/2$), and at Singapore ($\theta_0\approx 96°$): predict the daily rotation angle and discuss whether a Foucault pendulum would work in each location.
  \end{enumerate}
\end{newq}

\begin{newq}
  \NEW\; \textbf{Creative: Geodesics and Shortest Paths on Everyday Surfaces.}\;
  \begin{enumerate}[label=(\alph*)]
    \item A rubber band stretched around a cylinder follows a geodesic (helix). Compute its path explicitly and show the angle it makes with the axis is constant.
    \item A fly walks in a straight line on a cone (not through the apex). When the cone is unrolled, the path becomes a straight line. Is this path a geodesic? Justify using the isometry between the cone and a flat sector.
    \item A ship sailing from Singapore ($1.3^\circ$N, $103.8^\circ$E) to New York ($40.7^\circ$N, $74.0^\circ$W) follows a great-circle route. Describe qualitatively how the route looks on a Mercator map and why it appears curved.
  \end{enumerate}
\end{newq}

\begin{newq}
  \NEW\; \textbf{Creative: Turning Number and Geodesic Curvature.}\;
  The \emph{total geodesic curvature} of a closed curve $\gamma$ on $S$ is $\int_\gamma\kp_g\,ds$.
  \begin{enumerate}[label=(\alph*)]
    \item For a closed curve bounding a region $R$ on $S$, state the local Gauss--Bonnet formula: $\int_\gamma\kp_g\,ds + \iint_R K\,dA = 2\pi$.
    \item Apply it to a small circle of radius $\varepsilon$ on the unit sphere. Show $\int\kp_g\,ds\to 2\pi$ as $\varepsilon\to 0$ and $\int\kp_g\,ds < 2\pi$ for finite $\varepsilon$ (because $K=1>0$).
    \item Interpret: if $K>0$, closed curves ``turn less than in the plane''; if $K<0$, they ``turn more''. Relate this to the intuition that positively curved surfaces are ``bowl-shaped'' and negatively curved ones are ``saddle-shaped''.
  \end{enumerate}
\end{newq}

\begin{newq}
  \NEW\; \textbf{Creative: Index Theorem for Vector Fields.}\;
  The Poincar\'e--Hopf theorem states: for a smooth vector field on a compact surface with isolated zeros, $\sum_i\mathrm{ind}(p_i)=\chi(S)$.
  \begin{enumerate}[label=(\alph*)]
    \item Define the index of an isolated zero of a vector field.
    \item For $S^2$: the ``comb'' vector field (tangent to longitudes) has zeros only at the north and south poles, each with index $+1$. Verify $1+1=\chi(S^2)=2$.
    \item Explain why every smooth vector field on $S^2$ must have at least one zero (``you can't comb a hairy ball'') by using $\chi(S^2)=2\neq 0$.
  \end{enumerate}
\end{newq}

%----------------------------------------------------------
% Quick Reference Table
%----------------------------------------------------------
\newpage
\section*{Quick Reference: Key Formulas}
\addcontentsline{toc}{section}{Quick Reference: Key Formulas}

\begin{center}
\renewcommand{\arraystretch}{1.6}
\begin{tabular}{>{\bfseries}p{4.5cm} p{9cm}}
\toprule
Formula & Expression \\
\midrule
Curvature (non-unit speed) & $\kp = |\alpha'\times\alpha''|/|\alpha'|^3$ \\
Torsion (non-unit speed)   & $\tp = (\alpha'\times\alpha'')\cdot\alpha'''/|\alpha'\times\alpha''|^2$ \\
Arc length                 & $s(t) = \int_0^t |\alpha'(u)|\,du$, \quad $s'(t)=|\alpha'(t)|>0$ \\
Frenet matrix              & $(\mathbf{t}',\mathbf{n}',\mathbf{b}')^T = \begin{pmatrix}0&\kp&0\\-\kp&0&\tp\\0&-\tp&0\end{pmatrix}(\mathbf{t},\mathbf{n},\mathbf{b})^T$ \\
First fundamental form      & $\mathrm{I} = E\,du^2 + 2F\,du\,dv + G\,dv^2$ \\
Area element               & $dA = \sqrt{EG-F^2}\,du\,dv = |\mathbf{x}_u\wedge\mathbf{x}_v|\,du\,dv$ \\
Gaussian curvature          & $K = (eg-f^2)/(EG-F^2)$ \\
Mean curvature              & $H = (eG-2fF+gE)/\bigl(2(EG-F^2)\bigr)$ \\
Gauss map (graph)          & $N = (-f_x,-f_y,1)/\sqrt{1+f_x^2+f_y^2}$ \\
Normal curvature (Euler)   & $\kp_n = \kp_1\cos^2\theta + \kp_2\sin^2\theta$ \\
Gauss--Bonnet (global)      & $\iint_S K\,dA = 2\pi\chi(S)$ \\
Local Gauss--Bonnet         & $\int_{\partial R}\kp_g\,ds + \iint_R K\,dA + \sum_i\theta_i = 2\pi$ \\
Clairaut's relation         & $r\sin\psi = c$ (surface of revolution) \\
Sphere constraint           & $R^2 = 1/\kp^2 + (1/\tp)^2(1/\kp)'^2$ \\
Geodesic equations          & $\ddot{u}^k + \Gamma_{ij}^k\dot{u}^i\dot{u}^j = 0$ \\
Christoffel symbols        & $\Gamma_{ij}^k = \tfrac{1}{2}g^{kl}(\partial_i g_{jl}+\partial_j g_{il}-\partial_l g_{ij})$ \\
Jacobi equation            & $J''+K(\gamma)J=0$ along a geodesic \\
Covariant derivative        & $DV/dt = V' - \langle V',N\rangle N$ \\
\bottomrule
\end{tabular}
\end{center}

\end{document}
